\section{Panduan Visual NumPy}{Panduan Visual NumPy}
\label{appendix:numpy-visual-guide}
\objective{Operasi NumPy dapat sulit divisualisasikan, tetapi konsepnya sederhana.
Lampiran ini memberikan peragaan visual tentang penggunaan larik NumPy dengan sintaks pengirisan, penumpukan, penyiaran, dan operasi khusus sumbu.
Walaupun visualisasi ini ditujukan untuk larik berdimensi 1 atau 2, konsepnya dapat diperluas ke larik berdimensi $n$.
% See Lab \ref{lab:NumPy} for an introduction to NumPy operations and synatx.
}

\section*{Akses Data} % ======================================================

Entri-entri larik 2-D adalah baris-baris matriks tersebut (sebagai larik 1-D).
Untuk mengakses satu entri, masukkan indeks baris, tanda koma, lalu indeks kolom.
Ingat bahwa pengindeksan dimulai dari $0$.

\begin{align*}
\text{\li{A[0]}} = \left[\begin{array}{rrrrr}
\tikzmarkin{row1}\times & \times & \times & \times & \times\tikzmarkend{row1}\\
\times & \times & \times & \times & \times\\
\times & \times & \times & \times & \times\\
\times & \times & \times & \times & \times\\
\end{array}\right]
&&
\text{\li{A[2,1]}} = \left[\begin{array}{rrrrr}
\times & \times & \times & \times & \times\\
\times & \times & \times & \times & \times\\
\times & \tikzmarkin{entry} \times \tikzmarkend{entry} & \times & \times & \times\\
\times & \times & \times & \times & \times
\end{array}\right]
\end{align*}

\section*{Pengirisan} % =======================================================

Tanda titik dua tunggal mengekstrak seluruh baris atau kolom dari larik 2-D.
Sintaks \li{[a:b]} dapat dibaca sebagai ``entri ke-$a$ sampai (tetapi tidak termasuk) entri ke-$b$.''
Demikian pula, \li{[a:]} berarti ``entri ke-$a$ sampai akhir'' dan \li{[:b]} berarti ``semuanya sampai (tetapi tidak termasuk) entri ke-$b$.''


% The authority repeats this row/column example in two implementations.  Keep
% the later hf-tikz form, which carries the same mathematics and the intended
% highlights without the broken overlay's null-font semicolon artifacts.
\begin{align*}
\text{\li{A[1]}} = \text{\li{A[1,:]}} = \left[\begin{array}{rrrrr}
\times & \times & \times & \times & \times\\
\tikzmarkin{row2}\times & \times & \times & \times & \times\tikzmarkend{row2}\\
\times & \times & \times & \times & \times\\
\times & \times & \times & \times & \times\\
\end{array}\right]
&&
\text{\li{A[:,2]}} = \left[\begin{array}{rrrrr}
\times & \times & \tikzmarkin{col}\times & \times & \times\\
\times & \times & \times & \times & \times\\
\times & \times & \times & \times & \times\\
\times & \times & \times\tikzmarkend{col} & \times & \times
\end{array}\right]
\end{align*}

\begin{align*}
\text{\li{A[1:,:2]}} = \left[\begin{array}{rrrrr}
\times & \times & \times & \times & \times\\
\tikzmarkin{block}\times & \times & \times & \times & \times\\
\times & \times & \times & \times & \times\\
\times & \times\tikzmarkend{block} & \times & \times & \times
\end{array}\right]
&&
\text{\li{A[1:-1,1:-1]}} = \left[\begin{array}{rrrrr}
\times & \times & \times & \times & \times\\
\times & \tikzmarkin{interior} \times & \times & \times & \times\\
\times & \times & \times & \times \tikzmarkend{interior} & \times\\
\times & \times & \times & \times & \times\end{array}\right]
\end{align*}

\section*{Penumpukan} % =======================================================

\li{np.hstack()} menumpuk urutan larik secara horizontal, sedangkan \li{np.vstack()} menumpuk urutan larik secara vertikal.

\begin{align*}
\text{\li{A}} = \left[\begin{array}{ccc}
\textcolor{blue}{\times}&\textcolor{blue}{\times}&\textcolor{blue}{\times}\\
\textcolor{blue}{\times}&\textcolor{blue}{\times}&\textcolor{blue}{\times}\\
\textcolor{blue}{\times}&\textcolor{blue}{\times}&\textcolor{blue}{\times}
\end{array}\right]
&&
\text{\li{B}} = \left[\begin{array}{ccc}
\textcolor{red}{*} & \textcolor{red}{*} & \textcolor{red}{*} \\
\textcolor{red}{*} & \textcolor{red}{*} & \textcolor{red}{*} \\
\textcolor{red}{*} & \textcolor{red}{*} & \textcolor{red}{*}
\end{array}\right]
\end{align*}

\begin{align*}
\text{\li{np.hstack((A,B,A))}} =
\left[\begin{array}{ccccccccc}
\textcolor{blue}{\times}&\textcolor{blue}{\times}&\textcolor{blue}{\times}&
\textcolor{red}{*} & \textcolor{red}{*} & \textcolor{red}{*}&
\textcolor{blue}{\times}&\textcolor{blue}{\times}&\textcolor{blue}{\times}\\
\textcolor{blue}{\times}&\textcolor{blue}{\times}&\textcolor{blue}{\times}&
\textcolor{red}{*} & \textcolor{red}{*} & \textcolor{red}{*}&
\textcolor{blue}{\times}&\textcolor{blue}{\times}&\textcolor{blue}{\times}\\
\textcolor{blue}{\times}&\textcolor{blue}{\times}&\textcolor{blue}{\times}&
\textcolor{red}{*} & \textcolor{red}{*} & \textcolor{red}{*}&
\textcolor{blue}{\times}&\textcolor{blue}{\times}&\textcolor{blue}{\times}
\end{array}\right]
\end{align*}

\begin{align*}
\text{\li{np.vstack((A,B,A))}} =
\left[\begin{array}{ccc}
\textcolor{blue}{\times}&\textcolor{blue}{\times}&\textcolor{blue}{\times}\\
\textcolor{blue}{\times}&\textcolor{blue}{\times}&\textcolor{blue}{\times}\\
\textcolor{blue}{\times}&\textcolor{blue}{\times}&\textcolor{blue}{\times}\\
\textcolor{red}{*} & \textcolor{red}{*} & \textcolor{red}{*} \\
\textcolor{red}{*} & \textcolor{red}{*} & \textcolor{red}{*} \\
\textcolor{red}{*} & \textcolor{red}{*} & \textcolor{red}{*} \\
\textcolor{blue}{\times}&\textcolor{blue}{\times}&\textcolor{blue}{\times}\\
\textcolor{blue}{\times}&\textcolor{blue}{\times}&\textcolor{blue}{\times}\\
\textcolor{blue}{\times}&\textcolor{blue}{\times}&\textcolor{blue}{\times}
\end{array}\right]
\end{align*}
Karena larik 1-D berbentuk datar, \li{np.hstack()} menggabungkan larik-larik 1-D dan \li{np.vstack()} menumpuknya secara vertikal.
Untuk menjadikan beberapa larik 1-D sebagai kolom-kolom larik 2-D, gunakan \li{np.column_stack()}.

\begin{align*}
\text{\li{x}} = \left[\begin{array}{cccc}
\textcolor{blue}{\times}&\textcolor{blue}{\times}&\textcolor{blue}{\times}&\textcolor{blue}{\times}
\end{array}\right]
&&
\text{\li{y}} = \left[\begin{array}{cccc}
\textcolor{red}{*}&\textcolor{red}{*}&\textcolor{red}{*}&\textcolor{red}{*}\end{array}\right]
\end{align*}

\begin{align*}
\text{\li{np.hstack((x,y,x))}} =
\left[\begin{array}{cccccccccccc}
\textcolor{blue}{\times}&\textcolor{blue}{\times}&\textcolor{blue}{\times}&\textcolor{blue}{\times}&
\textcolor{red}{*}&\textcolor{red}{*}&\textcolor{red}{*}&\textcolor{red}{*}&
\textcolor{blue}{\times}&\textcolor{blue}{\times}&\textcolor{blue}{\times}&\textcolor{blue}{\times}
\end{array}\right]
\end{align*}

\begin{align*}
\text{\li{np.vstack((x,y,x))}} =
\left[\begin{array}{cccc}
\textcolor{blue}{\times}&\textcolor{blue}{\times}&\textcolor{blue}{\times}&\textcolor{blue}{\times}\\
\textcolor{red}{*}&\textcolor{red}{*}&\textcolor{red}{*}&\textcolor{red}{*}\\
\textcolor{blue}{\times}&\textcolor{blue}{\times}&\textcolor{blue}{\times}&\textcolor{blue}{\times}
\end{array}\right]
\end{align*}
\begin{align*}
\text{\li{np.column_stack((x,y,x))}} =
\left[\begin{array}{ccc}
\textcolor{blue}{\times}&\textcolor{red}{*}&\textcolor{blue}{\times}\\
\textcolor{blue}{\times}&\textcolor{red}{*}&\textcolor{blue}{\times}\\
\textcolor{blue}{\times}&\textcolor{red}{*}&\textcolor{blue}{\times}\\
\textcolor{blue}{\times}&\textcolor{red}{*}&\textcolor{blue}{\times}
\end{array}\right]
\end{align*}
Fungsi \li{np.concatenate()} dan \li{np.stack()} merupakan versi yang lebih umum dari \li{np.hstack()} dan \li{np.vstack()}, sedangkan \li{np.row_stack()} adalah alias untuk \li{np.vstack()}.

\section*{Penyiaran} % ========================================================

NumPy secara otomatis menyelaraskan larik untuk operasi per komponen kapan pun memungkinkan.
% The default behavior adds the first element to the first element of each row, the second element to the second element of each row, and so on.
Lihat \url{http://docs.scipy.org/doc/numpy/user/basics.broadcasting.html} untuk contoh yang lebih mendalam dan aturan penyiaran.

\begin{align*}
\text{\li{A}} = \left[\begin{array}{ccc}
1 & 2 & 3\\
1 & 2 & 3\\
1 & 2 & 3\\
\end{array}\right]
&&
\text{\li{x}} = \left[\begin{array}{ccc}
10 & 20 & 30
\end{array}\right]
\end{align*}

\begin{align*}
\text{\li{A + x}}
&= \begin{blockarray}{ccc}
\begin{block}{[ccc]}
1 & 2 & 3\\
1 & 2 & 3\\
1 & 2 & 3\\
\end{block}
  & + &  \\
\begin{block}{[ccc]}
10 & 20 & 30\\
\end{block}
\end{blockarray}
&= \left[\begin{array}{ccc}
11 & 22 & 33\\
11 & 22 & 33\\
11 & 22 & 33
\end{array}\right]
\\ \\
\text{\li{A + x.reshape((1,-1))}}
&= \left[\begin{array}{ccc}
1 & 2 & 3 \\
1 & 2 & 3 \\
1 & 2 & 3 \\
\end{array}\right]
+ \left[\begin{array}{c}
10 \\ 20 \\ 30\\
\end{array}\right]
&= \left[\begin{array}{ccc}
11 & 12 & 13\\
21 & 22 & 23\\
31 & 32 & 33\\
\end{array}\right]
\end{align*}

\section*{Operasi Sepanjang Sumbu} % ==========================================

Sebagian besar metode larik memiliki argumen \li{axis} yang memungkinkan suatu operasi dilakukan sepanjang sumbu tertentu.
Untuk menghitung jumlah setiap kolom, gunakan \li{axis=0}; untuk menghitung jumlah setiap baris, gunakan \li{axis=1}.

\begin{align*}
A = \left[\begin{array}{cccc}
1 & 2 & 3 & 4\\
1 & 2 & 3 & 4\\
1 & 2 & 3 & 4\\
1 & 2 & 3 & 4
\end{array}\right]
\end{align*}

\begin{align*}
\text{\li{A.<<sum>>(axis=0)}} &= %\text{\li{np.array([sum(A[:,j]) for j in range(A.shape[1])])}} =
\left[\begin{array}{cccc}
\tikzmarkin{col1}1 & \tikzmarkin{col2}2 & \tikzmarkin{col3}3 & \tikzmarkin{col4}4\\
1 & 2 & 3 & 4\\
1 & 2 & 3 & 4\\
1\tikzmarkend{col1} & 2\tikzmarkend{col2} & 3\tikzmarkend{col3} & 4\tikzmarkend{col4}
\end{array}\right]
= \left[\begin{array}{cccc} 4 & 8 & 12 & 16 \end{array}\right]
\\ \\
\text{\li{A.<<sum>>(axis=1)}} &= %\text{\li{np.array([sum(A[i,:]) for i in range(A.shape[0])])}} =
\left[\begin{array}{cccc}
\tikzmarkin{rowA}1 & 2 & 3 & 4\tikzmarkend{rowA}\\
\tikzmarkin{rowB}1 & 2 & 3 & 4\tikzmarkend{rowB}\\
\tikzmarkin{rowC}1 & 2 & 3 & 4\tikzmarkend{rowC}\\
\tikzmarkin{rowD}1 & 2 & 3 & 4\tikzmarkend{rowD}\\
\end{array}\right]
= \left[\begin{array}{cccc} 10 & 10 & 10 & 10 \end{array}\right]
\end{align*}
